\chapter*{Preface}
\addcontentsline{toc}{chapter}{Preface}

Book 1 closed by asking what the algebra of commutators looks like once it is asked to describe not merely its own structure but the curved manifold that structure implies. The present book takes up that question directly: it is the point at which the RSVP series stops treating geometry as a metaphor for algebraic curvature and starts building the geometry itself, with a connection, a metric, and a Ricci tensor recovered as derivatives of the same scalar-vector-entropy triple $\RSVPTriple$ that organized every construction in the algebraic cycle. Nothing here contradicts that cycle; everything here is what results from asking its commutator, quotient, and antifield structures to describe a manifold rather than merely an algebra.

The central substitution this book makes is thermodynamic curvature for spacetime curvature. Where ordinary differential geometry curves space by stretching it, the present book curves the plenum by smoothing it: the manifold's volume is held fixed throughout, $\det(g_{ij}) = \text{const}$, and whatever bending occurs is bending of the metric's internal structure under entropy flow rather than expansion of the underlying space. Gravitational and cognitive forces alike, on this account, are entropic geodesics, trajectories that minimize entropy production under this fixed-volume constraint, and this book provides the differential foundation for the entire geometric tier of the series: the bridge between the algebraic commutators of Cycle I and the categorical functors of curvature that Cycle III will develop.

